%% techbrief-cover.tex -- cover art for the TechBrief title banner.
%%
%% A vector redraw of the Issue 18 cover ("A" emblem in a broken ring with
%% connector nodes, over a circuit-board field).  Replace \tbcoverart (or
%% use \techbriefcover{...} in the document) to give an issue its own art.
%%
%% The picture is 202.9bp x 201.1bp.  Internally it is drawn in units of
%% the original 1158 x 1148 artwork, with y pointing up, so coordinates
%% can be read straight off a sketch.
\newdimen\tbpx \tbpx=0.17520bp % one pixel of the original artwork
\newcommand*\tbcoverart{%
\begin{tikzpicture}[x=0.17520bp,y=0.17518bp,line join=round,
    trace/.style={draw=tbctrace,line width=2.6\tbpx,line cap=round},
    pad/.style={draw=tbcpad,line width=2.6\tbpx,fill=tbcbgdark},
    via/.style={fill=tbctrace},
    glow/.style={fill=tbcglow}]
  \useasboundingbox (0,0) rectangle (1158,1148);
  \clip (0,0) rectangle (1158,1148);
  \coordinate (C) at (665,648);
  % ---------------------------------------------------------------
  % Background: deep navy, lifted behind the emblem.
  \fill[tbcbgdark] (0,0) rectangle (1158,1148);
  \shade[inner color=tbcbglight,outer color=tbcbgdark]
    (C) circle (760);
  % ---------------------------------------------------------------
  % Circuit traces: PCB-style runs that leave the ring radially,
  % break 45 degrees, and end on a pad.  Deterministic (fixed seed).
  \begin{scope}
    \clip (0,0) rectangle (1158,1148);
    \pgfmathsetseed{1810}
    \foreach \i in {1,...,96} {
      \pgfmathsetmacro\a{\i*3.75 + 1.8*rand}
      \pgfmathsetmacro\r{425 + 110*rnd}
      \pgfmathsetmacro\d{45*round(\a/45)}
      \pgfmathsetmacro\la{30 + 110*rnd}
      \pgfmathsetmacro\t{ifthenelse(rnd<0.5,-45,45)}
      \pgfmathsetmacro\lb{60 + 280*rnd}
      \pgfmathsetmacro\pad{rnd}
      \draw[trace] ($(C)+(\a:\r)$) coordinate (s)
        -- ++(\d:\la) -- ++(\d+\t:\lb) coordinate (e);
      \ifdim\pad pt<0.7pt \draw[pad] (e) circle (7.5);
      \else \fill[via] (e) circle (6); \fi
      \ifdim\pad pt<0.25pt \fill[via] (s) circle (5.5); \fi
    }
    % A second, sparser field toward the corners.
    \foreach \i in {1,...,72} {
      \pgfmathsetmacro\a{\i*5 + 2.5*rand}
      \pgfmathsetmacro\r{620 + 160*rnd}
      \pgfmathsetmacro\d{90*round(\a/90)}
      \pgfmathsetmacro\la{40 + 140*rnd}
      \pgfmathsetmacro\t{ifthenelse(rnd<0.5,-45,45)}
      \pgfmathsetmacro\lb{40 + 160*rnd}
      \draw[trace] ($(C)+(\a:\r)$) -- ++(\d:\la)
        -- ++(\d+\t:\lb) coordinate (e);
      \draw[pad] (e) circle (7.5);
    }
    % Glowing vias.
    \foreach \p in {(878,1048),(190,733),(192,693),(35,316),(340,223),
        (885,108),(1105,88),(1060,903),(880,218),(630,153),(95,1005),
        (1120,560)} {
      \fill[glow,opacity=.10] \p circle (22);
      \fill[glow,opacity=.22] \p circle (12);
      \fill[tbcglowcore] \p circle (6.5);
    }
  \end{scope}
  % ---------------------------------------------------------------
  % Emblem geometry.  Ring: centre C, radii 370..394.
  \def\tbarc##1##2{%
    ($(C)+(##1:394)$) arc[start angle=##1,end angle=##2,radius=394]
    -- ($(C)+(##2:370)$) arc[start angle=##2,end angle=##1,radius=370]
    -- cycle}
  % The "A": apex, tapered right leg with a rounded foot, and a left
  % leg that runs down into the ring and is trimmed by its outer edge.
  \def\tbA{%
    (600,828) [rounded corners=14\tbpx] -- (668,828) [rounded corners=0]
    -- (680,820) [rounded corners=8\tbpx] -- (892,404)
    [rounded corners=26\tbpx] -- (842,386) -- (790,443)
    [rounded corners=0] -- (640,693) -- (440,323) -- (300,300) -- cycle}
  \def\tbclipA{(C) circle (394)}
  % Connectors (centre lines) and nodes.
  \def\tbstemL{(360,638) .. controls (520,640) and (628,548) .. (623,340)}
  \def\tbstemR{(852,842) .. controls (826,786) and (782,744) .. (712,708)}
  % ---------------------------------------------------------------
  % Soft drop shadow for the whole emblem.
  \begin{scope}[shift={(10,-14)},tbcshadow]
    \fill \tbarc{167}{62};  \fill \tbarc{35}{-82};  \fill \tbarc{196}{236};
    \begin{scope}\clip \tbclipA; \fill \tbA; \end{scope}
    \draw[line width=26\tbpx] \tbstemL;  \draw[line width=22\tbpx] \tbstemR;
    \fill (315,638) circle (68);  \fill (623,277) circle (71);
    \fill (890,898) circle (72);
  \end{scope}
  % ---------------------------------------------------------------
  % Ring segments.
  \shade[left color=tbcsilver,right color=tbcsteel,shading angle=60]
    \tbarc{167}{62};
  \shade[top color=tbcblue,bottom color=tbcbluedeep] \tbarc{35}{-82};
  \shade[top color=tbcsilvermid,bottom color=tbcsteel] \tbarc{196}{236};
  % Right-hand connector runs under the A.
  \draw[line width=22\tbpx,tbcbluemid] \tbstemR;
  \shade[top color=tbcblue,bottom color=tbcbluedeep] (890,898) circle (72);
  % The A.
  \begin{scope}
    \clip \tbclipA;
    \shade[top color=tbcsilver,bottom color=tbcsteelblue,shading angle=-18]
      \tbA;
    \draw[white,opacity=.35,line width=2.5\tbpx] \tbA;
  \end{scope}
  % Left-hand connector crosses over the A's left leg; a dark keyline
  % cuts the leg so the connector reads as passing in front.
  \draw[line width=50\tbpx,tbcbgdark,opacity=.92] \tbstemL;
  \draw[line width=26\tbpx,tbcsilvermid] \tbstemL;
  \shade[top color=tbcsilver,bottom color=tbcsilvermid] (315,638) circle (68);
  \shade[top color=tbcsilvermid,bottom color=tbcsteel] (623,277) circle (71);
  % Bevel highlights on the nodes.
  \draw[white,opacity=.30,line width=2.5\tbpx] (315,638) circle (67);
  \draw[white,opacity=.25,line width=2.5\tbpx] (623,277) circle (70);
  \draw[white,opacity=.25,line width=2.5\tbpx] (890,898) circle (71);
\end{tikzpicture}}

% Palette (sampled from the original artwork).
\definecolor{tbcbgdark}{HTML}{050A15}
\definecolor{tbcbglight}{HTML}{0F1F3D}
\definecolor{tbctrace}{HTML}{18305E}
\definecolor{tbcpad}{HTML}{26457F}
\definecolor{tbcglow}{HTML}{3D74F0}
\definecolor{tbcglowcore}{HTML}{7FA6FF}
\definecolor{tbcsilver}{HTML}{CBD5E8}
\definecolor{tbcsilvermid}{HTML}{A2B1CE}
\definecolor{tbcsteel}{HTML}{7183B2}
\definecolor{tbcsteelblue}{HTML}{4A6BB4}
\definecolor{tbcblue}{HTML}{5F8BDB}
\definecolor{tbcbluemid}{HTML}{4F7ACC}
\definecolor{tbcbluedeep}{HTML}{3459AE}
\tikzset{tbcshadow/.style={black,opacity=.45}}
\endinput
